% Reusable content layouts. Copy the frame you need into your own section.
\begin{frame}{Compare two ideas}
  \begin{columns}[T,onlytextwidth]
    \begin{column}{.47\textwidth}
      \wvuheading{The question}
      \medskip

      What should the audience understand before seeing the result?
      \begin{itemize}
        \item State one concrete question.
        \item Define only the terms it needs.
        \item Keep the context brief.
      \end{itemize}
    \end{column}
    \begin{column}{.47\textwidth}
      \wvuheading{The evidence}
      \medskip

      What observation answers that question?
      \begin{itemize}
        \item Show the plot that answers it.
        \item Name its main limitation.
        \item End with the takeaway.
      \end{itemize}
    \end{column}
  \end{columns}
  \medskip
  \begin{information}[Layout tip]
    Top alignment with \texttt{[T]} keeps the two columns easy to compare.
  \end{information}
\end{frame}

\begin{frame}{Reveal an argument in three steps}
  \begin{enumerate}
    \item<1-> \textbf{Ask.} Begin with the question the audience needs answered.
    \item<2-> \textbf{Show.} Introduce the evidence that supports the result.
      \begin{itemize}
        \item<2-> Keep related detail with the claim it explains.
      \end{itemize}
    \item<3-> \textbf{Explain.} Connect the result to its practical meaning.
  \end{enumerate}
  \medskip
  \uncover<3->{\wvuheading{Each reveal adds one point.}}
  \note[item]<1>{Introduce the question before describing the method.}
  \note[item]<2>{Explain why this evidence answers the question.}
  \note[item]<3>{Pause on the takeaway before advancing.}
  % Presenter notes are exported by loading pgfpages and setting
  % \setbeameroption{show notes on second screen=right} in the preamble.
\end{frame}

\begin{frame}[label=technique-navigation]{Link to backup slides}
  Keep each main slide to \alert{one message}.
  Put longer explanations, settings, and backup material in the appendix.\footnote{A footnote holds a short note or a source.}

  \medskip
  \begin{beamercolorbox}[sep=.8em,wd=\linewidth]{palette primary}
    During questions, jump to the supporting detail and then return to
    this slide.
  \end{beamercolorbox}
  \medskip
  \hyperlink{reference-start}{\beamerskipbutton{Open the quick reference}}
\end{frame}
